\documentclass[a4paper, oneside]{article}
\usepackage{graphicx}
\usepackage{amsthm}
\usepackage{amsmath}
\usepackage{amssymb}
\usepackage[margin=0pt]{geometry}
\usepackage[italian]{babel}
\usepackage{pgfplots}
\usepackage{tabularx}
\usepackage{tikz-3dplot}
\usepackage{wrapfig}
\usepackage{color}
\usepackage{multicol}
\usepackage{arydshln}
\usepackage{mathtools}
\usepackage{enumerate}
\usepackage{graphicx}
\usepackage{svg}
\usepackage{cancel}
\usepackage[d]{esvect}
\usepackage[dvipsnames]{xcolor}
\usepackage{pifont}



%\usepackage{animate}
%\usepackage{xfp} % utile se vuoi fare calcoli aggiuntivi
\usetikzlibrary{tikzmark}
\newcommand{\TikzNCbar}[4][10pt]{
\tikz[overlay,remember picture]{\draw[#2] (#3) --++(0,-#1) -| (#4);}}
\newcommand\de{\partial}

\graphicspath{ {images/} }

\graphicspath{ {./images/} }
\usetikzlibrary{shapes.geometric}
\usetikzlibrary{datavisualization}
\usetikzlibrary{datavisualization.formats.functions}
\pgfplotsset{width=10cm,compat=1.9}

\setlength\dashlinedash{0.2pt}
\setlength\dashlinegap{1.5pt}
\setlength\arrayrulewidth{0.3pt}
\setlength{\multicolsep}{0pt}

\newcommand\eqq{\stackrel{\mathclap{\normalfont\mbox{?}}}{=}}
\newcommand\bulletout  {\labelitemfont \textbullet}
\newcommand\dvg{\operatorname{div}}
\newcommand\rot{\operatorname{rot}}
\newcommand{\tab}{\hspace*{2em}}
\newcommand{\xmark}{
\tikz[scale=0.23] {
    \draw[line width=0.7,line cap=round] (0,0) to [bend left=6] (1,1);
    \draw[line width=0.7,line cap=round] (0.2,0.95) to [bend right=3] (0.8,0.05);
}}
\newcommand{\cmark}{
\tikz[scale=0.23] {
    \draw[line width=0.7,line cap=round] (0.25,0) to [bend left=10] (1,1);
    \draw[line width=0.8,line cap=round] (0,0.35) to [bend right=1] (0.23,0);
}}
 \newcommand{\hookbox}[1]{
\begin{center}
\hfill\break
\begin{tikzpicture}
\node[inner sep=0pt,outer sep=0pt,anchor=base] (A) {
\begin{minipage}{\dimexpr\linewidth-5em}
\centering
#1
\end{minipage}
};
% Draw the left bracket
\draw ([xshift=0pt]A.north west) -- ++(0, 0.5) -- ++(0.4, 0);
% Draw the right bracket
\draw ([xshift=0pt]A.south east) -- ++(0, -0.5) -- ++(-0.4, 0);
\end{tikzpicture}
\end{center}}

\begin{document}

Data una superficie $\Sigma(u_1, u_2) = \begin{pmatrix}
    x(u_1, u_2) \\
    y(u_1, u_2) \\
    z(u_1, u_2)
\end{pmatrix}$ posso definire i suoi vettori coordinati come:
\begin{gather*}
    \frac{\de \Sigma (u_1, u_2)}{\de u_1} = \begin{pmatrix}
        \frac{\de x(u_1, u_2)}{\de u_1} \\
        \frac{\de y(u_1, u_2)}{\de u_1} \\
        \frac{\de z(u_1, u_2)}{\de u_1}
    \end{pmatrix}
    \qquad\qquad
    \frac{\de \Sigma (u_1, u_2)}{\de u_2} = \begin{pmatrix}
        \frac{\de x(u_1, u_2)}{\de u_2} \\
        \frac{\de y(u_1, u_2)}{\de u_2} \\
        \frac{\de z(u_1, u_2)}{\de u_2}
    \end{pmatrix}
\end{gather*}
Il vettore normale alla superficie è dato da:
\begin{gather*}
    \vv{N} = \frac{\de \Sigma (u_1, u_2)}{\de u_1} \times \frac{\de \Sigma (u_1, u_2)}{\de u_2}
\end{gather*}
ovvero:
\begin{gather*}
    \det\begin{pmatrix}
    \hat{i} & \hat{j} & \hat{k} \\
    \frac{\de x(u_1, u_2)}{\de u_1} & \frac{\de y(u_1, u_2)}{\de u_1} & \frac{\de z(u_1, u_2)}{\de u_1} \\
    \frac{\de x(u_1, u_2)}{\de u_2} & \frac{\de y(u_1, u_2)}{\de u_2} & \frac{\de z(u_1, u_2)}{\de u_2}
    \end{pmatrix}
\end{gather*}
E il modulo del vettore normale è dato da:
\begin{gather*}
    |\vv{N}| = \sqrt{N_x^2 + N_y^2 + N_z^2}
\end{gather*}
se la superficie è grafico di funzione ovvero:$\Sigma(u_1,u_2)=\begin{pmatrix}
    u_1 \\
    u_2 \\
    f(u_1, u_2)
\end{pmatrix}$ il vettore normale e la sua norma sono più veloci da calcolare e sono:
\begin{gather*}
    \vv{N} = \begin{pmatrix}
        -\frac{\de f(u_1, u_2)}{\de u_1} \\
        -\frac{\de f(u_1, u_2)}{\de u_2} \\
        1
    \end{pmatrix}
    \qquad\qquad
    |\vv{N}| = \sqrt{1 + \left(\frac{\de f(u_1, u_2)}{\de u_1}\right)^2 + \left(\frac{\de f(u_1, u_2)}{\de u_2}\right)^2}
\end{gather*}
Per calcolare il piano tangente alla superficie in un punto $P = (x_P, y_P, z_P)$ è sufficiente calcolare il piano che contiene i vettori coordinati della superficie in quel punto, ovvero:
\begin{gather*}
    \Pi_P = N_x(P)(x - x_P) + N_y(P)(y - y_P) + N_z(P)(z - z_P) = 0
\end{gather*}
Con $N_x(P), N_y(P), N_z(P)$ componenti del vettore normale alla superficie in $P$ che si calcolano in questo modo:
\begin{gather*}
    \begin{cases}
        N_x(u_1, u_2) = x_P\\
        N_y(u_1, u_2) = y_P\\
        N_z(u_1, u_2) = z_P
    \end{cases}
\end{gather*}


    Dato un corpo che occupa una regione di spazio $\Omega$, densità $\rho$ e distanza $\delta$ da un asse scelto: \\
    \bulletout L'area è $\int_{\Omega} 1 \ dx \ dy \ dz$ \\
    \bulletout La massa è $m=\int_{\Omega} \rho(x,y,z) \ dx \ dy \ dz$ \\
    \bulletout Il centro di massa $(B_x,B_y,B_z)$ è: $\begin{cases}
        B_x = \frac{1}{m} \int_{\Omega} x\cdot \rho(x,y,z) \ dx \ dy \ dz \\
        B_y = \frac{1}{m} \int_{\Omega} y\cdot \rho(x,y,z) \ dx \ dy \ dz \\
        B_z = \frac{1}{m} \int_{\Omega} z\cdot \rho(x,y,z) \ dx \ dy \ dz \\
    \end{cases}$ \\
    \bulletout Il momento di inerzia rispetto ad un generico asse $\tau$ è $M_\tau = \int_{\Omega} \delta_\tau^2 (x,y,z) \ \rho(x,y,z) \ dx \ dy \ dz$ e $\delta_\tau$ è la distanza da $\tau$

\end{document}